\chapter{What you need}
\label{ch:need}

\chapterbrief
  {Buy the right things, once.}
  {The decisions that are expensive to change later.}
  {A shopping list and no regrets.}

\section{The shopping list}

\begin{center}\small
\begin{tabular}{@{}L{0.30\linewidth}L{0.42\linewidth}r@{}}
\toprule
\thead{Item} & \thead{Notes} & \thead{Approx.} \\
\midrule
Raspberry Pi Zero 2 W & Any Pi works. See below. & \pounds18 \\
microSD card, 16\,GB+ & Class 10 / A1. A slow card makes everything feel broken. & \pounds8 \\
USB power supply, 5\,V 2.5\,A & \textbf{Not} a phone charger you had lying around. & \pounds8 \\
MFRC522 reader module & The ubiquitous red breakout board. & \pounds3 \\
\textbf{MIFARE DESFire EV2/EV3} & \textbf{Buy two. See the warning.} & \pounds3 ea \\
Female--female jumper wires & Get short ones, 10\,cm. Seriously. & \pounds3 \\
7-inch touchscreen \emph{(optional)} & 1024$\times$600, HDMI $+$ USB touch. Volume One, Chapter 12. & \pounds45 \\
\bottomrule
\end{tabular}
\end{center}

You also need a computer (macOS, Linux or Windows), a home network with Wi-Fi,
and an SD card reader.

\begin{trap}{The card matters more than everything else combined}
It must say \textbf{DESFire EV2} or \textbf{EV3}.

\textbf{Not} MIFARE Classic --- broken crypto, completely different protocol.
\textbf{Not} NTAG or Ultralight --- no crypto at all, no value files.
\textbf{Not} ``13.56\,MHz RFID tag'' from a bag of a hundred on an auction site.

If a listing says ``compatible'' or does not name the chip, it is not one.

\textbf{Buy at least two.} One to develop against, one spare --- because in
Volume One you will run a command that changes a key on the card, and a card
whose key you have changed and forgotten is scrap. Permanently, by design.
\end{trap}

\section{Which Pi?}

Any model with Wi-Fi works. The trade-off:

\begin{center}\small
\begin{tabular}{@{}L{0.24\linewidth}L{0.34\linewidth}L{0.34\linewidth}@{}}
\toprule
\thead{Model} & \thead{Good} & \thead{Annoying} \\
\midrule
Zero 2 W & Cheap, tiny, fits behind a screen & 512\,MB RAM; compiling on it is slow \\
Pi 3 / 4 / 5 & Fast, plenty of RAM & Bigger, more power, more money \\
Zero W (original) & Cheapest & Single core --- genuinely painful. Avoid. \\
\bottomrule
\end{tabular}
\end{center}

This book assumes a \textbf{Zero 2 W}, because it is the one that forces you to
do things properly: with 512\,MB you cannot be careless, and the workflow you
learn in Chapter~\ref{ch:deploy} --- build on the laptop, run on the Pi ---
is the right workflow on any hardware.

Everything works unchanged on a Pi 4; you will just wait less.

\begin{trap}{The power supply is not optional}
A Pi that browns out does not politely stop. It corrupts the SD card, drops the
network, and produces failures that look exactly like software bugs.

If you see random reboots, a rainbow square in the corner, or ``undervoltage
detected'' in \code{dmesg}, replace the supply before debugging anything else. A
phone charger rated ``5\,V 2\,A'' often will not hold 5\,V under load, and the
cable matters as much as the brick.
\end{trap}

\section{Wires: shorter than you think}

Get \textbf{10\,cm female-to-female} jumpers. The 20\,cm ones that come in most
kits are, at the 1\,MHz clock this project uses, genuinely unreliable --- and
their failure mode is intermittent, which is the worst kind.

\begin{analogy}{Shouting across a room}
A short wire is a conversation at arm's length. A long, unshielded, loosely
twisted jumper is shouting across a room with a fan running: most words arrive,
some do not, and you cannot predict which.

Digital signals do not degrade gracefully. A bit that arrives ambiguous is not
``slightly wrong'' --- it is a different number.
\end{analogy}

You will meet this again in Appendix~\ref{app:trouble0}: \emph{version register
reads a different value each run} is almost always wire length, and almost never
code.

\section{What you do not need}

\begin{itemize}
  \item \textbf{A monitor, keyboard or mouse for the Pi.} This book sets it up
        \emph{headless} --- configured before first boot, reached over the
        network. It is faster, and it is how real deployments work.
  \item \textbf{A soldering iron.} The RC522 usually ships with header pins
        already attached. If yours came with loose pins in a bag, you do need to
        solder them, or buy one that is pre-soldered.
  \item \textbf{An Arduino, a level shifter, or a breadboard.} The Pi's GPIO is
        already 3.3\,V, which is what the RC522 wants.
\end{itemize}

\begin{checkpoint}{Before continuing}
You have, physically in front of you: a Pi, an SD card, a power supply, an
RC522, jumper wires, and at least one card that says \textbf{DESFire EV2} or
\textbf{EV3} on the packaging.

If the card is the only thing missing, carry on --- you will not need it until
Chapter~\ref{ch:wiring}.
\end{checkpoint}
